\documentclass[11pt,a4paper]{article}
\usepackage[margin=1in]{geometry}
\usepackage{amsmath,amssymb,amsthm}
\usepackage{algorithm}
\usepackage{algpseudocode}
\usepackage{booktabs}
\usepackage{hyperref}

\theoremstyle{definition}
\newtheorem{definition}{\textbf{Definition}}[section]
\newtheorem{theorem}{\textbf{Theorem}}[section]
\newtheorem{lemma}[theorem]{\textbf{Lemma}}
\newtheorem{remark}{\textbf{Remark}}[section]

\title{\textbf{Spiking Neural Networks and Leaky Integrate-and-Fire Dynamics}}
\author{\textbf{Team Scaibu} \\ \small Email: \texttt{jaydeep.9664920749@gmail.com} \\ \small \textbf{Architecture and Core Engine Team}}
\date{\textbf{October 2026}}

\begin{document}
\maketitle

\section{Definitions and Cognitive Mental Models}

\subsection{Formal Mathematical Definition}
\textbf{Definition (Continuous Leaky Integration, Discontinuous Spike Generation, and Surrogate Derivatives):}
In computational neuroscience and neuromorphic engineering, the \textbf{Leaky Integrate-and-Fire (LIF)} model describes the subthreshold membrane voltage $V(t)$ of an electrically polarized neuron governed by the continuous-time linear differential equation:
{\boldmath
\begin{equation}
\tau_m \frac{\mathrm{d}V(t)}{\mathrm{d}t} = - (V(t) - V_{\text{rest}}) + R_m I(t)
\end{equation}
}
where $\tau_m = R_m C_m$ is the membrane time constant, $V_{\text{rest}}$ is the resting equilibrium potential (conveniently set to 0), and $I(t)$ is the incoming synaptic current.

In discrete-time simulation with step size $\Delta t$, the subthreshold state updates via the exponential decay factor $\beta = \exp(-\Delta t / \tau_m) \in (0, 1)$:
{\boldmath
\begin{equation}
V[t] = \beta V[t-1] + I[t]
\end{equation}
}
A spike $S[t] \in \{0, 1\}$ is emitted instantaneously when the membrane potential reaches or surpasses the threshold $V_{\text{th}} > 0$:
{\boldmath
\begin{equation}
S[t] = \Theta(V[t] - V_{\text{th}}) = 
\begin{cases}
1, & \text{if } V[t] \ge V_{\text{th}} \\
0, & \text{if } V[t] < V_{\text{th}}
\end{cases}
\end{equation}
}
Following a spike emission ($S[t] = 1$), the membrane voltage resets immediately to a baseline state according to either:
{\boldmath
\begin{align}
\text{Hard Reset:} \quad V[t] &\gets 0 \\
\text{Soft Reset:} \quad V[t] &\gets V[t] - V_{\text{th}}
\end{align}
}
Because the Heaviside step function $\Theta$ has zero derivative almost everywhere and undefined derivative at $0$, standard backpropagation through time (BPTT) fails. Spiking neural networks resolve this using a smooth \textbf{Surrogate Gradient} $\sigma'(x)$, such as the fast arc-tangent derivative:
{\boldmath
\begin{equation}
\frac{\partial S}{\partial V} \approx \frac{1}{\pi} \frac{\alpha}{1 + \alpha^2 (V - V_{\text{th}})^2}
\end{equation}
}
where $\alpha > 0$ controls the surrogate gradient bandwidth.

\subsection{Expanded Non-Technical Definition (Intuition for Anyone)}
\textbf{Intuitive Conceptual Definition:}
\begin{enumerate}
    \item \textbf{How the Human Brain Sips 20 Watts of Power}: Modern GPU clusters consume megawatts of power running giant continuous matrix multiplications 100\% of the time, even when nothing changes. The human brain computes at extreme efficiency because billions of neurons stay completely silent until a critical spark of electricity occurs.
    \item \textbf{Binary Spikes Instead of Floating-Point Decimals}: In standard AI, every layer outputs dense 32-bit floating numbers ($0.8741, -2.1932$). In spiking neural networks, neurons communicate using binary Morse-code pulses: $1$ (spike) or $0$ (silence).
    \item \textbf{Multiplications Turn into Additions}: When a neuron is silent ($0$), downstream synapses perform zero operations. When a neuron fires ($1$), the receiver simply adds the weight into its bucket without needing expensive multiply-accumulate (MAC) circuitry!
    \item \textbf{Surrogate Gradients (Training the Un-differentiable)}: You cannot calculate standard derivatives across an on/off light switch. Surrogate gradients replace the step-function switch with a gentle virtual ramp during the backward pass so the AI can learn with standard gradient descent.
\end{enumerate}

\subsection{Child-Friendly Mental Model (Physical Metaphor)}
Imagine a leaky water bucket under a dripping faucet:
\begin{itemize}
    \item \textbf{The Leaky Hole}: The bucket has a small hole at the bottom. If water drips slowly, it leaks right out and the bucket stays empty ($\beta V$).
    \item \textbf{The Sudden Rush of Rain}: If a sudden torrent of rain pours in ($I[t]$), the water level rises faster than it can leak out.
    \item \textbf{The Tipping Latch}: When the water hits the red line at the top ($V_{\text{th}}$), the bucket trips a mechanical bell: \emph{DING!} (Spike $S=1$!).
    \item \textbf{The Dump and Reset}: The bell trips a trapdoor that instantly empties the bucket back to zero, ready to collect fresh raindrops!
\end{itemize}

\section{Core Mathematical Equations and Definitions}

\subsection{Subthreshold Discrete Leaky Integration}
{\boldmath
\begin{equation}
V_{\text{decayed}}[t, i] = \beta \cdot V[t-1, i] + I[t, i]
\end{equation}
}
\begin{itemize}
    \item \textbf{Formal Definition}: The first-order autoregressive decay and accumulation of input synaptic current.
    \item \textbf{What It Means}: Blends past membrane charge with fresh external electrical charge while leaking energy exponentially.
    \item \textbf{Why It Is Important}: Gives each neuron temporal memory, allowing it to integrate information arriving asynchronously over time.
\end{itemize}

\subsection{Discontinuous Heaviside Firing Condition}
{\boldmath
\begin{equation}
S[t, i] = \Theta(V_{\text{decayed}}[t, i] - V_{\text{th}})
\end{equation}
}
\begin{itemize}
    \item \textbf{Formal Definition}: The step-threshold operator producing binary events $\{0, 1\}$.
    \item \textbf{What It Means}: Emits a binary communication spike if the membrane potential surpasses the action potential threshold.
    \item \textbf{Why It Is Important}: Enforces event-driven temporal sparsity, cutting energy consumption on neuromorphic hardware.
\end{itemize}

\subsection{Membrane Reset Dynamics (Hard vs Soft)}
{\boldmath
\begin{equation}
V[t, i] = 
\begin{cases}
V_{\text{decayed}}[t, i] \cdot (1 - S[t, i]), & \text{Hard Reset} \\
V_{\text{decayed}}[t, i] - S[t, i] \cdot V_{\text{th}}, & \text{Soft Reset}
\end{cases}
\end{equation}
}
\begin{itemize}
    \item \textbf{Formal Definition}: The state reset mechanism discharging the membrane potential following a spike.
    \item \textbf{What It Means}: Hard reset clears the membrane to 0; soft reset preserves residual excess charge above threshold.
    \item \textbf{Why It Is Important}: Soft reset prevents information loss caused by throwing away surplus excitation energy.
\end{itemize}

\subsection{Fast Arc-Tangent Surrogate Gradient Operator}
{\boldmath
\begin{equation}
\sigma'(u) = \frac{1}{\pi} \frac{\alpha}{1 + \alpha^2 u^2}, \quad u = V_{\text{decayed}} - V_{\text{th}}
\end{equation}
}
\begin{itemize}
    \item \textbf{Formal Definition}: The continuous, bell-shaped Cauchy-Lorentz surrogate derivative substituting the Dirac delta distribution.
    \item \textbf{What It Means}: Provides smooth, non-zero gradient flow around the threshold boundary during backpropagation.
    \item \textbf{Why It Is Important}: Resolves the vanishing/dead gradient problem of discontinuous spiking neurons.
\end{itemize}

\section{Rigorous Mathematical Analysis Checklist}

\begin{itemize}
    \item \textbf{Notation}: $V$ (membrane voltage), $I$ (synaptic current), $\beta$ (leak factor), $V_{\text{th}}$ (firing threshold), $S$ (binary spike train), $\sigma'$ (surrogate gradient).
    \item \textbf{Epistemic Status}: Classical bio-physics formulation (Lapicque 1907), modernized for deep learning via surrogate gradients by Neftci et al. (2019).
    \item \textbf{Dependency Graph}: $V[t-1] \to V_{\text{decayed}}[t] \to (S[t], \sigma'[t]) \to V[t] \to V[t+1]$.
    \item \textbf{Dimensions}: Currents $I \in \mathbb{R}^{T \times N}$, spikes $S \in \{0, 1\}^{T \times N}$, voltage $V \in \mathbb{R}^N$.
    \item \textbf{Regimes}: Neuromorphic event-driven computing (Intel Loihi, SynSense Speck), ultra-low-power edge audio/vision processing.
    \item \textbf{Invariants}: Spike states $S[t, i] \in \{0, 1\}$; subthreshold leak strictly dissipates energy: $\beta \in [0, 1)$.
    \item \textbf{Isomorphisms}: LIF with surrogate gradients is isomorphic to a Recurrent Neural Network (RNN) with quantized step activation and identity state decay.
\end{itemize}

\section{Theoretical Theorems and Formal Proofs}

\begin{theorem}[Asymptotic Steady-State Voltage under Constant DC Current]
Let a subthreshold LIF neuron receive constant input current $I[t] = I_0$ for all $t \ge 0$, with leak factor $\beta \in [0, 1)$ and initial voltage $V[0] = 0$. If $I_0 < (1 - \beta) V_{\text{th}}$, the neuron will never fire, and its membrane potential converges asymptotically to the steady-state limit:
{\boldmath
\begin{equation}
\lim_{t \to \infty} V[t] = \frac{I_0}{1 - \beta}
\end{equation}
}
\end{theorem}

\begin{proof}
Since $I_0 < (1 - \beta) V_{\text{th}}$, assume no spikes are emitted ($S[t] = 0$).
The discrete linear recurrence is:
{\boldmath
\begin{equation}
V[t] = \beta V[t-1] + I_0
\end{equation}
}
Unrolling the recurrence from $V[0] = 0$:
{\boldmath
\begin{equation}
V[t] = I_0 \sum_{k=0}^{t-1} \beta^k = I_0 \frac{1 - \beta^t}{1 - \beta}
\end{equation}
}
Since $\beta \in [0, 1)$, $\lim_{t \to \infty} \beta^t = 0$.
Taking the limit as $t \to \infty$:
{\boldmath
\begin{equation}
\lim_{t \to \infty} V[t] = \frac{I_0}{1 - \beta}
\end{equation}
}
Since by assumption $I_0 < (1 - \beta) V_{\text{th}}$, we have $\frac{I_0}{1 - \beta} < V_{\text{th}}$.
Because $V[t]$ is strictly monotonically increasing and upper-bounded by $\frac{I_0}{1 - \beta} < V_{\text{th}}$, the membrane potential never crosses the threshold $V_{\text{th}}$, verifying that no spikes are emitted.
\end{proof}

\begin{theorem}[Bounded Mean Firing Rate]
Let $T$ denote the total simulation horizon. The temporal mean firing rate $\bar{r}_i = \frac{1}{T} \sum_{t=0}^{T-1} S[t, i]$ of any LIF neuron is strictly bounded on the unit interval:
{\boldmath
\begin{equation}
0 \le \bar{r}_i \le 1
\end{equation}
}
\end{theorem}

\begin{proof}
At every discrete time step $t \in \{0, \dots, T-1\}$, the spike emission is defined by the indicator function $S[t, i] \in \{0, 1\}$.
Summing over all $T$ time steps:
{\boldmath
\begin{equation}
0 \le \sum_{t=0}^{T-1} S[t, i] \le \sum_{t=0}^{T-1} 1 = T
\end{equation}
}
Dividing the entire inequality by $T > 0$:
{\boldmath
\begin{equation}
0 \le \frac{1}{T} \sum_{t=0}^{T-1} S[t, i] \le 1
\end{equation}
}
Hence $0 \le \bar{r}_i \le 1$.
The lower bound is attained when $I[t, i] < (1-\beta)V_{\text{th}}$ (quiescent neuron), and the upper bound is attained when $I[t, i] \ge V_{\text{th}}$ for all $t$ under hard reset (maximum firing saturation).
\end{proof}

\section{Foundational Architectural Questions and Epistemic Meta-Questions}

\subsection{Architectural Question 1: Information Transmission via Rate vs Temporal Coding}
\textbf{Question:} How does an SNN encode information in spike trains?\\
\textbf{Answer:} There are two primary paradigms: (1) \emph{Rate Coding}, where information is proportional to the total spike count $\bar{r}$ over window $T$, requiring multiple time steps to convey high precision; and (2) \emph{Temporal/Time-to-First-Spike Coding}, where the exact microsecond timestamp of the first spike conveys information with single-spike energy efficiency.

\textbf{Meta-Question:} Which coding scheme is most compatible with surrogate gradient BPTT?\\
\textbf{Answer:} Rate-like dense burst coding trains reliably with standard Adam optimizers because surrogate gradients flow across many active time steps. Sparse single-spike temporal coding suffers from severe gradient vanishing unless exact spike-timing loss functions (like EventProp or SpikeGrad) are used.

\subsection{Architectural Question 2: Hard Reset vs Soft Reset}
\textbf{Question:} Why does soft reset outperform hard reset in deep SNN architectures?\\
\textbf{Answer:} Under hard reset, if a neuron accumulates $V = 1.9$ against threshold $V_{\text{th}} = 1.0$, the voltage resets to $0$, discarding $0.9$ units of excitation charge. Under soft reset, the voltage resets to $1.9 - 1.0 = 0.9$, carrying over the residual energy. This guarantees exact charge conservation and eliminates quantization distortion.

\textbf{Meta-Question:} Is soft reset biologically realistic?\\
\textbf{Answer:} Biological potassium and sodium ion channels force a hyperpolarization refractory period closer to hard reset. However, in artificial deep networks, soft reset acts as a high-fidelity continuous integrator that maximizes signal-to-noise ratio.

\section{Algorithmic Implementation Specification}

\begin{algorithm}[H]
\caption{Leaky Integrate-and-Fire Dynamic Simulation with Surrogate Gradients}
\begin{algorithmic}[1]
\Require Currents $\mathbf{I} \in \mathbb{R}^{T \times N}$, initial voltage $\mathbf{V}_0 \in \mathbb{R}^N$, leak $\beta \in [0, 1]$, threshold $V_{\text{th}} > 0$, reset mode $\in \{\text{hard}, \text{soft}\}$, surrogate scale $\alpha > 0$
\Ensure Spike trains $\mathbf{S} \in \{0, 1\}^{T \times N}$, final voltage $\mathbf{V}_{\text{final}} \in \mathbb{R}^N$, firing rates $\mathbf{r} \in \mathbb{R}^N$, surrogate grads $\mathbf{G} \in \mathbb{R}^{T \times N}$
\State $\mathbf{V} \gets \operatorname{copy}(\mathbf{V}_0), \quad \text{counts} \gets \mathbf{0} \in \mathbb{Z}^N$
\State Initialize $\mathbf{S} \in \mathbb{Z}^{T \times N}, \quad \mathbf{G} \in \mathbb{R}^{T \times N}$
\For{$t = 0$ \textbf{to} $T-1$}
    \For{$i = 0$ \textbf{to} $N-1$}
        \State $v_{\text{dec}} \gets \beta \cdot \mathbf{V}[i] + \mathbf{I}[t][i]$ \Comment{Leaky Integration}
        \State $diff \gets v_{\text{dec}} - V_{\text{th}}$
        \State $spike \gets (diff \ge 0.0) \;?\; 1 : 0$ \Comment{Threshold Check}
        \State $u \gets \alpha \cdot diff$ \Comment{Arc-Tangent Surrogate Derivative}
        \State $grad \gets \frac{1}{\pi} \cdot \frac{\alpha}{1.0 + u^2}$
        \If{$spike = 1$} \Comment{State Reset}
            \State $\text{counts}[i] \gets \text{counts}[i] + 1$
            \If{reset mode is soft}
                \State $\mathbf{V}[i] \gets v_{\text{dec}} - V_{\text{th}}$
            \Else
                \State $\mathbf{V}[i] \gets 0.0$
            \EndIf
        \Else
            \State $\mathbf{V}[i] \gets v_{\text{dec}}$
        \EndIf
        \State $\mathbf{S}[t][i] \gets spike, \quad \mathbf{G}[t][i] \gets grad$
    \EndFor
\EndFor
\State Initialize $\mathbf{r} \in \mathbb{R}^N$
\For{$i = 0$ \textbf{to} $N-1$}
    \State $\mathbf{r}[i] \gets \text{counts}[i] / T$
\EndFor
\State \Return $(\mathbf{S}, \mathbf{V}, \mathbf{r}, \mathbf{G})$
\end{algorithmic}
\end{algorithm}

\section{Big-$\Theta$ Complexity Analysis}

\begin{itemize}
    \item \textbf{Dynamic Integration Time Complexity}: $\Theta(T \cdot N)$ scalar operations across all simulation time steps and neurons.
    \item \textbf{Surrogate Derivative Evaluation}: $\Theta(T \cdot N)$ operations.
    \item \textbf{Space Complexity}: $\Theta(T \cdot N)$ memory buffers for output spike trains and surrogate computational tapes.
    \item \textbf{Neuromorphic Hardware Operations}: If total spikes emitted is $M \ll T \cdot N$, synaptic compute complexity drops to $\Theta(M \cdot d_{\text{synapse}})$ sparse additions with zero multiplications.
    \item \textbf{Parallel Depth}: $\mathcal{D} = \Theta(T)$ sequential recurrence steps in time; $\Theta(1)$ parallel across the $N$ neurons.
\end{itemize}

\section{Single Canonical Repository Reference}

The complete deterministic scalar implementation, verification suite, and unit test benchmarks are published at:
\begin{center}
\url{https://github.com/Chief-Strategist-J/llm-obs-infra}
\end{center}

\end{document}
